\documentclass[12pt]{article}
\usepackage[utf8]{inputenc}
\usepackage[T1]{fontenc}
\usepackage{hyperref}
\title{Tsakhur Language Resources for K-12 Curriculum Development}
\author{Compiled by Claude Sonnet 5 (claude-sonnet-5) for Andr\'as Kornai}
\date{}
\begin{document}
\maketitle

\section{Introduction}
Tsakhur (ISO 639-3 tkr) is a Nakh-Daghestanian (Lezgic) language spoken in northwestern Azerbaijan and southern Dagestan, Russia. Materials use Latin script on the Azerbaijani side and Cyrillic in Dagestan.

\section{Primary Educational Resources}
Wikipedia is reported to say that Tsakhur is taught in primary schools in both countries. This conflicts with the project's reference description and was not verified against a primary source. A textbook by Jafarov was mentioned in one result, with no title, date or URL. No textbook URL was found.

\section{Linguistic and Reference Materials}
\begin{itemize}
\item Tsakhur -- Azerbaijani -- Russian -- English Dictionary (SIL, on Webonary, about 3,743 Tsakhur entries): \url{https://www.webonary.work/tsakhur/?p=146817}
\item Freelang Tsakhur--English dictionary (about 600 entries each way): \url{https://www.freelang.net/dictionary/tsakhur.php}
\item Wolfgang Schulze, \emph{Tsakhur} (Lincom, 1997), grammar sketch; Glottolog record: \url{https://glottolog.org/resource/reference/id/119500}
\item Omniglot page on Tsakhur (alphabet and sample text; cites Ibragimov, \emph{Tsakhurskij jazyk}, Moscow 1990): \url{https://omniglot.com/writing/tsakhur.htm}
\item Tsakhur language, Wikipedia (mentions Dirr's 1913 grammar): \url{https://en.wikipedia.org/wiki/Tsakhur_language}
\item Manfred Ostrowski's academia.edu page, where a review of a large Tsakhur grammar was seen: \url{https://independent.academia.edu/OstrowskiManfred}
\end{itemize}

\section{Religious and Community Resources}
None found. A general page about the people and language was returned (\url{https://shakasana.com/en/node/51}) and was not verified.

\section{Digital and Audio Resources}
The Webonary dictionary above is online. No audio or video was found.

\section{Curriculum Development Context}
The language is spoken by all adults and by most young people. This search was in English only. Russian- and Azerbaijani-language searches should turn up school textbooks and local publications, and these may exist. A curriculum would need to handle the Latin and Cyrillic orthographies separately.
\end{document}